\writeSectionFrame{\presentationTitle}{Fazit}
\begin{frame}{Vorteile}
	\begin{itemize}
	  \item Klare Trennung von Regeln und Ablauf
	  \item Hohe Wiederverwendbarkeit
	  \item Sehr gut testbar
	  \item Lesbarer Code
	\end{itemize}
\end{frame}

\begin{frame}{Nachteile}
	\begin{itemize}
	  \item Mehr Klassen
	  \item Kann für einfache Regeln überdimensioniert wirken
	  \item Initial höherer Implementierungsaufwand
	\end{itemize}
\end{frame}

\begin{frame}{Specification und Testbarkeit}
	\begin{itemize}
	  \item Jede Specification ist isoliert testbar
	  \item Keine Abhängigkeit von UI oder Infrastruktur
	  \item Keine komplexen Test-Setups nötig
	\end{itemize}
\end{frame}

\begin{frame}{Was Specification nicht ist}
	\begin{itemize}
	  \item Kein Ersatz für:
	  \begin{itemize}
	    \item komplexe Workflows
	    \item Algorithmen
	    \item Persistenzlogik
	  \end{itemize}
	  \item Fokus auf:
	  \begin{itemize}
	    \item Entscheidungslogik
	    \item Geschäftsregeln
	  \end{itemize}
	\end{itemize}
\end{frame}

\begin{frame}{Zusammenspiel mit anderen Mustern}
	\begin{itemize}
	  \item Dependency Injection
	  \item Fluent API
	  \item Immutable Objects
	  \item Domain-Driven Design
	\end{itemize}
\end{frame}

\begin{frame}{Vorteile der Fluent-Verpackung}
	\begin{itemize}
	  \item Höhere Lesbarkeit: fast Sätze in natürlicher Sprache
	  \item Intuitive Bedienung
	  \item Kettenbildung
	  \item Weniger syntaktischer Lärm: keine verschachtelten new-Aufrufe mehr
	\end{itemize}
\end{frame}

\begin{frame}{Fazit: Specification Pattern}
	\begin{itemize}
	  \item Macht Geschäftsregeln explizit
	  \item Fördert sprechende Namen statt if-Kaskaden
	  \item Sehr gut kombinierbar mit:
	  \begin{itemize}
	    \item Dependency Injection
	    \item Fluent APIs
	    \item Immutable Objects
	  \end{itemize}
	  \item Mehr Designarbeit – aber weniger Chaos
	\end{itemize}
\end{frame}
